\documentclass[11pt]{article}
\usepackage[a4paper,margin=25mm]{geometry}
\usepackage{amsmath,amssymb,amsthm}
\usepackage[hidelinks]{hyperref}
\newtheorem{theorem}{Theorem}
\setlength{\emergencystretch}{1em}
\begin{document}
\title{Incompatible clauses in Conjecture 2171}
\author{Chengzhe Gao}\date{}\maketitle
The source asserts both that Frankl's union-closed conjecture holds and
that the ground-set bound twelve for a minimal counterexample is attained.
These two clauses are incompatible. This observation neither proves nor
disproves the classical Frankl conjecture considered alone.

For a finite set family $\mathcal F$, put
$d_{\mathcal F}(x)=|\{A\in\mathcal F:x\in A\}|$.
Use the ground set $\bigcup\mathcal F$, and exclude the degenerate family
whose only possible member is the empty set. The usual Frankl assertion is
\[
 \forall\mathcal F\text{ union-closed},\quad
 \exists x\in\bigcup\mathcal F:\quad
 2d_{\mathcal F}(x)\geq|\mathcal F|.
\]
A counterexample is an admissible union-closed family for which
$2d_{\mathcal F}(x)<|\mathcal F|$ for every ground element.

\begin{theorem}
Frankl's assertion and the existence of a minimal counterexample with
ground-set size twelve cannot both hold.
\end{theorem}
\begin{proof}
Assume both statements, and take the claimed family on twelve ground
elements. Frankl's assertion applies to that family and supplies an element
$x$ with $2d_{\mathcal F}(x)\geq|\mathcal F|$. The counterexample condition
gives the strict opposite inequality for the same $x$, a contradiction.
Minimality and the value twelve are not needed for the contradiction.
\end{proof}

The words ``with twelve attained'' in English and the corresponding Chinese phrase (``and attained'')
are an existential assertion, not merely a conditional lower bound on
hypothetical counterexamples. Removing attainment would remove this
contradiction and would be a different problem.

\paragraph{Formal verification.}
\texttt{Family n} is a list of actual subsets of $\mathrm{Fin}(n)$, represented
by Boolean membership functions. \texttt{Admissible} requires no repeated
sets, closure under pointwise union, full ground-set support, and the
presence of a nonempty set. Every finite set family with an $n$-element
ground set admits such a representation by labeling its ground elements;
union and occurrence counts are preserved. Conversely each admissible
representation is such a family. \texttt{frequency} computes the actual
number of sets containing an element. The theorem
\texttt{counterexample\_\allowbreak{}iff} connects the strict inequalities to absence
of a Frankl witness. \texttt{MinimalGroundCounterexample} requires an
actual counterexample and the absence of any smaller-ground counterexample.
The final \texttt{conjecture2171\_\allowbreak{}false} negates the conjunction of the
universal Frankl assertion and an actual attained twelve-element minimum.
Thus the formalization does not replace the set theory by an uninterpreted
proposition. It proves only the inconsistency of the source's full statement.

Lean 4.19.0 with \texttt{Std} checks the proof, with only standard logical
axioms in the printed dependency lists.

\end{document}
